\newpage
\section{Missió 6: Empipa dos planetes sencers}\label{sec:mission6}

\begin{minipage}{0.87\linewidth}
Han passat anys a la colònia des que vas ajudar a configurar la gran antena per comunicar-vos amb la Terra. Fa poc, la colònia ha descobert \textit{una altra} colònia secreta, a Saturn -- i la Terra ni es va molestar a dir-nos-ho!\\

\noindent
L'assemblea ha decidit fer una petita broma i et torna a demanar ajuda.
El teu objectiu és segrestar la propera connexió TCP Saturn $\leftrightarrow$ Terra
i després enviar un memoràndum fals fent-te passar per Saturn. L'assemblea demana que sigui elegant.
\end{minipage}
%
\hspace{0.03\linewidth}
%
\begin{minipage}{0.10\linewidth}
\includegraphics[width=\linewidth]{ostrich}
\end{minipage}

\subsection*{Detalls de l'informe}
\noindent
L'enginyera en cap de la colònia ha fet una exploració i està disposada a ajudar-te.
Ha recopilat la informació següent:

\begin{itemize}[itemsep=0.1cm]
\item La teva HugeAntenna\texttrademark\ de la colònia es comunica amb una altra HugeAntenna\texttrademark\ de la Terra. La Terra fa servir aquesta mateixa antena per comunicar-se amb una altra HugeAntenna\texttrademark\ de Saturn.
Quan la Terra rep el primer paquet d'una connexió TCP, en verifica l'origen.

\item Després d'aquest primer paquet, la Terra no verifica l'origen de cap paquet. De fet, els missatges falsos o inesperats posteriors s'ignoren, sense activar cap alarma. Calculant bé el moment, ho aprofitaràs per injectar els teus paquets falsos.

\item El primer dia de cada mes, exactament a les 4:20 (hora de Nova Barcelona, Terra), la Terra rep el primer paquet d'una connexió TCP de Saturn. Aquestes connexions serveixen per intercanviar memoràndums oficials, i l'assemblea vol que en segrestis la propera.

\item Saturn usa una pila TCP/IP obsoleta i insegura amb aquest defecte: sempre posa el número de seqüència TCP a $0$ al primer missatge d'una connexió TCP.

\item No podràs veure els missatges entre la Terra i Saturn. Però l'enginyera en cap va capturar aquesta frame (trama) Ethernet, que va iniciar la connexió TCP anterior (diguem-ne $C_{-1}$) entre Saturn i la Terra (\texttt{00 d0} va seguit de \texttt{f1 bb} i \texttt{45 10} va seguit de \texttt{00 28}):

\vspace{0.25cm}
\begin{center}
\begin{minipage}{0.5\linewidth}
\begin{verbatim}
00 09 bf 42 07 64 00 d0   f1 bb 0f 5a 08 00 45 10
00 28 d2 71 40 00 40 06   00 00 ca f7 44 05 4c 4c
15 15 9f ee 00 a9 00 00   00 00 00 00 00 00 50 02
01 00 00 00 00 00
\end{verbatim}
\end{minipage}
\end{center}
\vspace{0.25cm}
\end{itemize}

\noindent
Amb això, l'assemblea demana un informe amb els aspectes següents:

\begin{enumerate}[itemsep=0.1cm]
\item Dibuixa un diagrama swimlane de l'establiment de la connexió anterior ($C_{-1}$) entre Saturn i la Terra.
A més, tabula \textit{tots} els camps de capçalera d'aquests paquets (i els seus valors).

\item Explica el teu pla d'atac, \textit{és a dir}, com segrestaràs la connexió. Indica clarament:
    \begin{itemize}
    \item Quines suposicions fas sobre la propera connexió (p. ex., quins camps seran idèntics o podrien canviar).
    \item Quants missatges enviaràs i quan tens previst enviar-los.
    \end{itemize}

\item Mostra la seqüència completa de bytes d'una frame que contingui el teu memoràndum de broma.
\end{enumerate}

\subsection*{Detalls del retorn}

Segueix un procés de retorn com el de les missions anteriors, tractant aquests aspectes clau:
\begin{itemize}
\item S'ha descodificat bé la frame capturada per l'enginyera en cap?
\item Has ideat un pla que necessiti menys missatges?
\item Has pogut descodificar el memoràndum de broma? Era elegant?
\end{itemize}

